\documentclass[11pt,a4paper]{article}
\usepackage[utf8]{inputenc}
\usepackage[T1]{fontenc}
\usepackage[spanish,provide=*]{babel}
\addto\captionsspanish{\renewcommand{\tablename}{Tabla}}
\usepackage{lmodern}
\usepackage{microtype}
\usepackage[margin=2.5cm]{geometry}
\usepackage{amsmath,amssymb}
\usepackage{graphicx}
\usepackage{booktabs}
\usepackage{adjustbox}
\usepackage{pdflscape}
\usepackage{float}
\usepackage[font=small,labelfont=bf]{caption}
\usepackage{enumitem}
\usepackage{xcolor}
\usepackage{listings}
\usepackage[round]{natbib}
\usepackage[hidelinks]{hyperref}
\setlength{\parskip}{0.5em}
\setlength{\parindent}{0pt}
\setlist{itemsep=0.2em}

% Texto en rojo = algo que hay que revisar/actualizar después de volver a correr el do-file.
% Para sacar el color: \renewcommand{\revisar}[1]{#1}
\newcommand{\revisar}[1]{\textcolor{red}{#1}}
% Inserta una tabla exportada desde Stata (se achica si no entra)
\newcommand{\tabla}[1]{\begin{adjustbox}{max width=\linewidth, max totalheight=0.72\textheight}\input{#1}\end{adjustbox}}
% Figura que todavía no existe (se genera al correr el do-file): muestra un recuadro
\newcommand{\figura}[2][0.75\linewidth]{\IfFileExists{#2}{\includegraphics[width=#1]{#2}}{\fbox{\parbox{\dimexpr#1-2\fboxsep-2\fboxrule\relax}{\centering\vspace{2em}\textcolor{red}{Figura pendiente: se genera al correr el do-file}\\[0.5em]\texttt{\detokenize{#2}}\vspace{2em}}}}}
% Notas al pie de tablas
\newcommand{\notas}[1]{\par\vspace{0.3em}\begin{minipage}{0.92\linewidth}\footnotesize\textit{Notas:} #1\end{minipage}}
% Consignas en cursiva
\newcommand{\consigna}[1]{\par\textit{#1}\par}

\lstset{
  language={},
  basicstyle=\ttfamily\scriptsize,
  breaklines=true,
  columns=fullflexible,
  keepspaces=true,
  showstringspaces=false,
  frame=single,
  rulecolor=\color{gray!50},
  commentstyle=\color{gray},
  morecomment=[l]{*\ },
  morecomment=[l]{//},
  morecomment=[s]{/*}{*/},
  tabsize=4,
  inputencoding=utf8,
  extendedchars=true,
  literate={á}{{\'a}}1 {é}{{\'e}}1 {í}{{\'i}}1 {ó}{{\'o}}1 {ú}{{\'u}}1
           {Á}{{\'A}}1 {É}{{\'E}}1 {Í}{{\'I}}1 {Ó}{{\'O}}1 {Ú}{{\'U}}1
           {ñ}{{\~n}}1 {Ñ}{{\~N}}1 {ü}{{\"u}}1 {Ü}{{\"U}}1
           {¿}{{?`}}1 {¡}{{!`}}1 {’}{{'}}1 {‘}{{`}}1 {“}{{``}}1 {”}{{''}}1
           {ã}{{\~a}}1 {õ}{{\~o}}1 {â}{{\^a}}1 {ê}{{\^e}}1 {ô}{{\^o}}1 {à}{{\`a}}1
           {Ã}{{\~A}}1 {Õ}{{\~O}}1 {Â}{{\^A}}1 {Ê}{{\^E}}1 {Ô}{{\^O}}1 {Ç}{{\c C}}1 {×}{{$\times$}}1
           {°}{{$^\circ$}}1 {–}{{--}}1 {—}{{---}}1 {·}{{$\cdot$}}1 {ç}{{\c c}}1
}
\begin{document}

\begin{center}
{\Large\bfseries Problem Set 4: Variables instrumentales}\\[0.5em]
E178: Economía Aplicada --- Universidad de San Andrés\\
Fecha de entrega: 13 de septiembre de 2024\\[0.5em]
\textbf{Integrantes:} Juan Batalla, Nina Di Costanzo Pereira, Catalina Gallardo e Ignacio Yenarópulos
\end{center}

Este trabajo se basa en \citet{murphy}, que usa la presencia de inmigrantes chinos a principios del siglo XX como instrumento de la presencia actual de cárteles de droga en los municipios de México, para estimar el efecto de los cárteles sobre distintos resultados socioeconómicos. Usamos la base \texttt{poppy.dta}, que no contiene todas las variables del paper; en particular, no tenemos la medida principal de presencia de cárteles (2011), así que trabajamos con la presencia en 2010.

\section*{1. Variable \textit{Chinese presence}}

Siguiendo al paper, \textit{Chinese presence} es una variable binaria que vale 1 si en el municipio vivía al menos una persona nacida en China en 1930 (\texttt{chinos1930hoy} $>0$) y 0 si no vivía ninguna. La variable tiene 16 valores faltantes, que corresponden a las delegaciones del Distrito Federal.

\section*{2. Estadística descriptiva}

\begin{table}[H]\centering
\caption{Estadísticas descriptivas}\label{tab:desc}
\small
\tabla{"../output/Tabla 1.tex"}
\notas{Datos a nivel municipal; incluye al Distrito Federal. Cada variable se resume con todas sus observaciones no faltantes. \textit{Marginalization} es el índice de marginación de 2015 (valores más altos indican mayor marginación). \textit{Chinese population} es la cantidad de personas nacidas en China que residían en el municipio en 1930.}
\end{table}

La Tabla~\ref{tab:desc} presenta las estadísticas descriptivas. En 2010 había presencia de cárteles en el 29\% de los municipios, mientras que en 2005 solo en el 9\%, lo que refleja la expansión de los cárteles durante la guerra contra el narcotráfico. El 14\% de los municipios tenía población china en 1930, pero esa población estaba muy concentrada: el promedio es de 4.2 personas por municipio, con un máximo de 1,183. El índice de marginación está estandarizado (media 0 y desvío 1). Entre los controles, el 10\% de los municipios tenía población alemana en 1930 (variable que el paper usa como placebo) y el 9\% tiene un clima apto para el cultivo de amapola. Los valores son muy parecidos a los de la Tabla 1 del paper. Las variables de población en 1930 y de crecimiento poblacional tienen alrededor de 2,160 observaciones, lo que define el tamaño de la muestra en las especificaciones con controles. De aquí en adelante excluimos al Distrito Federal.

\section*{3. Regresiones por MCO (Tabla 5, columnas 3 a 6)}

\begin{table}[H]\centering
\caption{Presencia china y presencia de cárteles (MCO)}\label{tab:ols}
\small
\tabla{"../output/Tabla 2.tex"}
\notas{Errores estándar clusterizados por estado entre paréntesis. Todas las regresiones incluyen efectos fijos por estado. Los controles son: German presence, Poppy suitability, temperatura y precipitación promedio, superficie, población en 1930, distancias a EE.UU., a la Ciudad de México y al puerto más cercano, y Head of state. El wild bootstrap p-valor surge de 999 repeticiones. * $p<0.10$, ** $p<0.05$, *** $p<0.01$.}
\end{table}

Al incluir efectos fijos por estado, como hace el paper, replicamos exactamente las columnas 3 a 6 de su Tabla 5 (Tabla~\ref{tab:ols}). Los municipios con presencia china en 1930 tienen una probabilidad de presencia de cárteles en 2010 21.6 puntos porcentuales mayor que los municipios del mismo estado sin presencia china. Al controlar por características geográficas y demográficas, la diferencia cae a 10.1 puntos porcentuales, pero sigue siendo significativa (p-valor de wild bootstrap de 0.014). Para 2005 los efectos son menores (13.0 y 5.9 puntos porcentuales), en línea con la menor presencia de cárteles en ese año, y también son significativos (p-valores de wild bootstrap de 0.000 y 0.015). Usamos el wild cluster bootstrap porque hay solo 30 o 31 clusters, y con pocos clusters los errores clusterizados convencionales tienden a subestimar la incertidumbre. La columna 2 es la primera etapa de la especificación básica del punto siguiente: muestra que el instrumento es relevante, aunque su poder explicativo cae bastante al agregar controles.

\section*{4. Estimaciones por variables instrumentales (Tablas 7 y 8)}

\begin{table}[H]\centering
\caption{Presencia de cárteles y marginación (2SLS)}\label{tab:iv}
\small
\tabla{"../output/Tabla 3.tex"}
\notas{Variable dependiente: índice de marginación 2015. \textit{Cartel presence} (2010) está instrumentada con \textit{Chinese presence}. Errores estándar clusterizados por estado entre paréntesis. El F de la primera etapa corresponde al test del instrumento excluido. Columna (1): sin controles. Columna (2): especificación básica. Columna (3): excluye los municipios a menos de 100 km de la frontera con EE.UU. Columna (4): excluye Sinaloa. Columna (5): agrega el crecimiento poblacional local 1920--30. El wild bootstrap p-valor surge de 999 repeticiones. * $p<0.10$, ** $p<0.05$, *** $p<0.01$.}
\end{table}

El estimador de variables instrumentales usa solo la parte de la variación en la presencia de cárteles que se explica por la presencia histórica de inmigrantes chinos, que (según el argumento del paper) introdujeron el cultivo de amapola y dieron origen a las primeras redes de tráfico. Con un solo instrumento, el estimador es el cociente entre la forma reducida (efecto de la presencia china sobre la marginación, $-0.241$ en la especificación básica) y la primera etapa (efecto de la presencia china sobre la presencia de cárteles, $0.101$): $-0.241/0.101\approx-2.39$.

La Tabla~\ref{tab:iv} muestra que la presencia de cárteles reduce el índice de marginación en alrededor de 2 a 2.8 puntos, es decir, entre 2 y 2.8 desvíos estándar, y el efecto es significativo en todas las especificaciones, también con los p-valores de wild bootstrap (entre 0.000 y 0.017). En la especificación básica (columna 2) el coeficiente es $-2.39$, idéntico al que reporta el paper para la presencia de cárteles en 2010 (Tabla 9, columna 2). La Tabla 7 del paper usa la presencia de cárteles en 2011 y obtiene coeficientes menores en valor absoluto ($-1.90$ en la especificación básica): como la forma reducida es la misma, el estimador de IV cambia en proporción inversa a la primera etapa. El resultado es robusto a excluir los municipios de la frontera ($-2.80$), a excluir Sinaloa ($-2.42$) y a controlar por el crecimiento poblacional de 1920--30 ($-2.38$). Por MCO, en cambio, el coeficiente es mucho menor ($-0.39$). Hay que notar dos cosas. Primero, el efecto es enorme: llevaría a un municipio desde la media hasta cerca del mínimo de la distribución. Segundo, el F de la primera etapa es bajo cuando se incluyen controles (7.7 en la especificación básica y 5.6 en la columna 3, por debajo de la regla práctica de 10), lo que plantea un problema de instrumento débil.

\begin{table}[H]\centering
\caption{Presencia de cárteles y componentes del índice de marginación (2SLS)}\label{tab:comp}
\small
\tabla{"../output/Tabla 4.tex"}
\notas{Variables dependientes en porcentaje de la población (o de las viviendas) de cada municipio. Especificación básica: controles y efectos fijos por estado. \textit{Cartel presence} (2010) está instrumentada con \textit{Chinese presence}. Errores estándar clusterizados por estado entre paréntesis. El wild bootstrap p-valor surge de 999 repeticiones. * $p<0.10$, ** $p<0.05$, *** $p<0.01$.}
\end{table}

La Tabla~\ref{tab:comp} descompone el efecto en los componentes del índice. Según los p-valores de wild bootstrap, la presencia de cárteles reduce significativamente al 5\% el analfabetismo ($-21.8$ puntos porcentuales), la población sin primaria completa ($-34.3$ p.p.), la proporción de viviendas sin electricidad ($-4.6$ p.p.), la población en localidades pequeñas ($-100.7$ p.p.) y la población con ingresos bajos ($-27.2$ p.p.). El efecto sobre el hacinamiento ($-20.0$ p.p.) es significativo al 10\% (p-valor de 0.051). Los efectos sobre las viviendas sin baño, sin agua y con piso de tierra no son significativos. El patrón es el mismo que el de la Tabla 8 del paper (nuestros coeficientes son un 25\% más grandes por usar la presencia de cárteles en 2010). La interpretación intuitiva es que, en los municipios donde la presencia china histórica llevó a una mayor presencia de cárteles, los cárteles habrían generado actividad económica, empleo e ingresos, y financiado bienes públicos locales, lo que se refleja en mejores indicadores de educación, vivienda e ingresos. Sin embargo, las magnitudes son poco creíbles: un efecto de $-100.7$ puntos porcentuales sobre la población en localidades pequeñas equivale a pasar de un municipio totalmente rural a uno totalmente urbano. Esto refuerza la cautela sobre el instrumento (ver puntos 5 y 6).

\section*{5. Tests de exogeneidad}

\textbf{Exogeneidad del instrumento.} No es posible testearla. Tenemos un solo instrumento para una sola variable endógena, de modo que el modelo está exactamente identificado y no hay restricciones de sobreidentificación que testear: los tests de Sargan o Hansen requieren más instrumentos que variables endógenas. De hecho, por construcción los residuos de 2SLS son ortogonales al instrumento, así que regresar los residuos contra el instrumento siempre da un coeficiente nulo y no aporta información. La exogeneidad del instrumento es un supuesto que solo se puede defender con argumentos y con evidencia indirecta (por ejemplo, tests de placebo como los que hace el paper con la presencia alemana o con los condados de EE.UU.).

\textbf{Exogeneidad de la presencia de cárteles.} Si supiéramos que el instrumento es exógeno, sí podríamos testear si la presencia de cárteles es endógena con el test de Hausman (Durbin--Wu--Hausman). Bajo la hipótesis nula de que la variable es exógena, tanto MCO como 2SLS son consistentes (y MCO es más eficiente), por lo que sus estimaciones no deberían diferir sistemáticamente. Bajo la alternativa, solo 2SLS es consistente. Como nuestros errores están clusterizados por estado, usamos la versión del test basada en una regresión auxiliar, que es robusta (\texttt{estat endogenous}). El test equivale a agregar el residuo de la primera etapa a la ecuación estructural y testear si su coeficiente es cero. En la especificación básica obtenemos $F(1,29)=10.03$, con un p-valor de 0.004, así que rechazamos la hipótesis nula al 1\%: la presencia de cárteles es endógena, MCO es inconsistente y, si el instrumento es válido, debe preferirse la estimación por variables instrumentales. El test de Hausman clásico (comando \texttt{hausman}) no es válido con errores clusterizados, por eso no lo usamos.

\section*{6. Crítica}

El supuesto crítico del paper es la restricción de exclusión: la presencia china en 1930 solo debería afectar la marginación actual a través de la presencia de cárteles. Esto es difícil de sostener. Los inmigrantes chinos no se distribuyeron al azar: se instalaron en lugares con oportunidades económicas (el noroeste, cerca de la frontera, de los puertos del Pacífico, del ferrocarril y de la minería) y muchos se dedicaron al comercio. Esas características, y la propia actividad comercial de las comunidades chinas, pueden haber afectado directamente el desarrollo de largo plazo de esos municipios. Los controles por distancias y población en 1930 capturan solo una parte de estas diferencias, y la variación histórica que no capturan termina atribuyéndose a los cárteles.

Además, el instrumento es débil una vez que se incluyen controles (F de 7.7 en la especificación básica). Con un instrumento débil, cualquier violación de la exclusión se amplifica: el sesgo asintótico del estimador de IV es $\text{cov}(Z,u)/\text{cov}(Z,X)$, de modo que con una primera etapa de alrededor de 0.10, un efecto directo pequeño de la presencia china sobre la marginación se multiplica por diez. Las magnitudes implausibles de las Tablas~\ref{tab:iv} y~\ref{tab:comp} (más de dos desvíos estándar en el índice, más de 100 puntos porcentuales en localidades pequeñas) son consistentes con este problema. Con instrumentos débiles, además, la inferencia convencional no es confiable, por lo que convendría reportar intervalos de Anderson--Rubin.

Por último, el canal propuesto (inmigración china, cultivo de amapola, cárteles) sugiere otras vías por las que el instrumento puede afectar los resultados: la economía de la droga en su conjunto (cultivo, empleo agrícola, lavado de dinero) puede influir en el desarrollo local aunque no se manifieste como presencia de un cártel según la medida basada en noticias de Google, que además tiene error de medición. Los ejercicios de placebo del paper (presencia alemana, condados de EE.UU.) son útiles pero indirectos, porque esas poblaciones y esos lugares difieren en aspectos importantes. Una forma de evaluar la sensibilidad de los resultados sería acotar el efecto bajo violaciones moderadas de la exclusión (por ejemplo, con el enfoque de ``instrumentos plausiblemente exógenos'' de Conley, Hansen y Rossi).

\begin{thebibliography}{9}
\bibitem[Murphy y Rossi(2020)]{murphy} Murphy, T. E., y Rossi, M. A. (2020). Following the poppy trail: Origins and consequences of Mexican drug cartels. \textit{Journal of Development Economics}, 143, 102433.
\end{thebibliography}

\newpage
\appendix
\section*{Apéndice: código}
\lstinputlisting{../programs/ps4.do}

\end{document}
